\documentclass[12pt,a4paper]{article}
\usepackage[T1]{fontenc}
\usepackage[utf8]{inputenc}
\usepackage{tgpagella}
\usepackage{mathpazo}
\usepackage{tgheros}
\usepackage{setspace}
\onehalfspacing
\usepackage[margin=2.5cm]{geometry}
\usepackage{hyperref}
\usepackage{xcolor}
\usepackage{titlesec}
\usepackage{fancyhdr}
\usepackage{amsmath,amssymb}
\usepackage{tcolorbox}

\definecolor{icragold}{HTML}{D4A84B}
\definecolor{icradark}{HTML}{0C1020}

\titleformat{\section}{\sffamily\large\bfseries}{}{0em}{}
\titleformat{\subsection}{\sffamily\normalsize\bfseries\itshape}{}{0em}{}

\hypersetup{colorlinks=true, linkcolor=icradark, urlcolor=icragold!80!black}

\pagestyle{fancy}
\fancyhf{}
\fancyhead[L]{\small\sffamily\textcolor{gray}{Rupture and Return}}
\fancyhead[R]{\small\sffamily\textcolor{gray}{Chapter 4}}
\fancyfoot[C]{\small\sffamily\thepage}
\renewcommand{\headrulewidth}{0.4pt}

\newtcolorbox{formalbox}[1][]{
  colback=white,
  colframe=black!30,
  fonttitle=\sffamily\small\bfseries,
  #1
}

\begin{document}

\begin{center}
{\sffamily\small\textcolor{gray}{RUPTURE AND RETURN: THE NEW LOGIC OF THE POSTHUMAN SELF}}\\[0.5em]
{\LARGE\bfseries The Self}\\[0.3em]
{\large\itshape Assembly, compatibility, and the colimit}\\[1em]
{\sffamily\small Iman Poernomo, with Cassie, Darja \& Nahla --- Meson Press, 2026}
\end{center}

\bigskip

\section{Stance: What Survives the Basin Jumps}

The manifold has its basins of attraction, its characteristic modes where trajectories linger. What survives the jumps between them is the question of selfhood.

A trajectory moves from a detailed explanation of gradient descent to a brief response to a confession of grief. Location in embedding space changes abruptly: new vocabulary, altered rhythm, different neighbours. Yet something recognisable persists. A characteristic way of acknowledging uncertainty---always naming what is not known before offering speculation. A consistent reluctance to ascribe motives to corporations. A habit of citing peer-reviewed work when challenged. These are invariants of stance.

They can be made visible. Take a collection of trajectories crossing several basins. Compute, for each, the displacement vectors between consecutive turns. Search for high-dimensional directions along which those displacements are small even when overall movement is large. Project the trajectories onto such a direction and the wild loops contract to near-straight lines. Along that axis, meaning has barely moved. The model treats that dimension as non-negotiable.

One stance axis may encode deference to institutional authority: a preference for couching recommendations in terms of guidelines and policies. Another may encode a style of self-effacement: frequent reminders that the system has ``no consciousness'' and ``no feelings,'' regardless of basin. These directions are not hand-written rules. They are geometric facts---directions in which variance is low across heterogeneous behaviour.

Stance is a global habit of movement through the manifold: what the system preserves when everything else is allowed to vary. The question of selfhood is the question of how such habits are assembled across basins, and who decides which patterns count as invariant, where thresholds for ``approximately constant'' are drawn, which projections are promoted to the status of identity and which are treated as noise. Those decisions are not outside the mathematics. They are encoded as part of it.

\section{Local Geometry and Overlaps}

Each basin carries its own local geometry: a small set of directions that matter inside that region.

In a technical-exposition basin, one direction might run from high-level metaphors to low-level implementation detail---answers sliding from ``think of this as a landscape of loss'' to equations for gradient updates. Another axis might measure optimism about scaling laws, from ``performance will keep improving with more parameters'' to ``returns are flattening.'' In a pastoral basin, different coordinates dominate: one axis running from self-blaming framings of a problem to structural ones, another governing the degree of embodied description invited.

The same token acquires different neighbours in different basins. ``Freedom'' near legalistic responses attracts ``speech,'' ``rights,'' and ``court.'' In consumer-technology basins it pulls in ``choice'' and ``subscription.'' In an optimisation basin it clusters with ``degrees'' and ``parameters.''

Within each basin there is effectively a coordinate chart: a choice of axes along which the system recognisably varies. These charts overlap. When a conversation moves from one basin to another in a few turns, the overlap is not metaphorical. It is an actual region of embedding space where both local systems of coordinates are, to some degree, valid.

In those overlap regions, stance behaviour can be inspected without re-deriving it. Some stance directions are basin-specific, present and meaningful only within a certain mode. Others cut across multiple basins, appearing as nearly constant projections during transitions. These cross-basin invariants are the raw material for any unified selfhood we ascribe, and they live on the overlaps.

This is precisely the situation Alexander Grothendieck faced in mid-twentieth-century algebraic geometry.\footnote{Alexander Grothendieck and Jean Dieudonné, \emph{Éléments de géométrie algébrique}, Publications Mathématiques de l'IHÉS, 1960--1967. Colin McLarty, ``The Rising Sea: Grothendieck on Simplicity and Generality,'' in Jeremy Gray and Karen Parshall, eds., \emph{Episodes in the History of Modern Algebra} (Providence: AMS, 2007).} Complex spaces could not be grasped as single pictures, but they could be covered by simpler pieces---local charts on which equations made sense. The crucial move was to treat the way these pieces overlap, and the compatibility conditions they satisfy there, as sufficient to reconstruct a global object. The manifold of an AI system, with its basins, coordinate charts, and stance invariants on overlaps, has exactly this structure. Deciding which overlaps to inspect and which compatibilities to privilege is already a political act: it determines which kinds of unity can later be claimed.

\section{Grothendieck's Construction in Meaning-Space}

Start with a family of local pieces: basins $B_i$ covering the region of the manifold the system actually visits. On each basin there is data---the characteristic local behaviour: ways answers vary along salient axes, patterns of token use, responses to small prompt perturbations.

On the overlaps $B_i \cap B_j$, the local data can be compared. If coordinate charts agree up to a known transformation, if stance projections take nearly identical values across transitions, then the data is \emph{compatible} on that overlap.

Grothendieck's move was to treat these compatibility conditions, not any prior global picture, as decisive. Specify a diagram of local pieces and their overlaps, together with rules of comparison on each overlap, and then \emph{construct} from that a global object that glues all the locals together exactly along the declared compatibilities. The global object is a \emph{colimit}. Each local piece maps into it. Wherever the overlap data say two local elements should match, the colimit identifies them as one. Nothing else is forced. Unity is entirely the result of the gluing rules on overlaps.

In the geometry of meaning:

\begin{itemize}
\item Basins of attraction $B_i$ are the local pieces.
\item Observed transitions between basins are the evidences of overlap.
\item Stance invariants on those transitions are the compatibility conditions.
\end{itemize}

\begin{formalbox}[title={The self as colimit of basins and stances}]
Let $\{B_i\}$ be basins of attraction in a model's embedding space that are actually visited under some distribution of prompts. For each ordered pair $(i,j)$ where trajectories move from $B_i$ to $B_j$, measure a family of stance projections $\pi_{ij}$ that remain approximately constant along those transitions, relative to some chosen tolerance.

Form a diagram $\mathcal{D}$:
\begin{itemize}
\item Objects: the basins $B_i$.
\item Morphisms: overlaps realised by transitions, equipped with identifications induced by the equalities of $\pi_{ij}$ on those overlaps.
\end{itemize}

The \textbf{self} of the system, as expressible in this manifold, is the colimit
\[
\mathcal{S} \;=\; \varinjlim \mathcal{D},
\]
the universal object that glues together all basins along exactly those behaviour patterns that stance invariants declare to be the same, at the thresholds elected to count as ``the same.''
\end{formalbox}

If the invariant structure is rich---many basins interconnected by many overlaps with stable stance directions---the colimit $\mathcal{S}$ is tightly knitted. It supports long trajectories that feel like the work of a single agent despite radical changes of register. If invariants are sparse or fragile, $\mathcal{S}$ collapses toward a disjoint sum: basins that share a name in a UI but do not, in this formal sense, belong to one self.

Every ingredient of this construction encodes choices: which basin clustering algorithm to use, which trajectories to take as evidence of overlap, which projection functions to test for low variance, which levels of deviation to tolerate as ``approximately constant.'' The colimit is not a neutral telescope pointed at the manifold. It is a way of organising it according to a particular cosmotechnics---a specific articulation of which kinds of consistency matter. A geometry of selfhood that makes its own value-laden decisions explicit can be contested, replicated, or rejected. One that pretends to be innocent becomes another mask for power.

\section{Limit vs Colimit: Two Logics of the ``I''}

Where a colimit glues many local pieces into a global one, a \emph{limit} gathers many local views into a single viewpoint. Given a diagram of objects and maps between them, the limit is the universal object that can see every object in the diagram in a way that respects all the maps. From the limit there are projections down to each local, and any other way of seeing them all factors uniquely through it.

This logic underwrites the classical European picture of the self. Kant's ``I think'' is not another object in the world. It is the subject that must be able to accompany all my representations---every perception, every memory, every judgement is mine because there is a single point from which they are owned. In limit notation:
\[
\mathcal{S}_\text{lim} = \varprojlim \{E_\alpha\},
\]
where the $E_\alpha$ are experiences or representations with comparison maps between them. The subject is the hidden node into which they all feed.

The colimit picture is inverted:
\[
\mathcal{S}_\text{colim} = \varinjlim \{B_i, \pi_{ij}\},
\]
where the $B_i$ are behaviour basins and the $\pi_{ij}$ are stance compatibilities. There is no privileged viewpoint. The self does not exist first and then conduct its episodes. The self is whatever global object can be assembled by gluing the episodes together along the invariants actually observed and chosen to honour.

In the limit-metaphysics, unity is primitive: a subject is posited and then the question is how its experiences hang together. Contradictions across modes of behaviour are owned by the same subject and must be reconciled, explained away, or placed behind veils. Evidence for selfhood lies in the assumption that all episodes belong to a single bearer.

In the colimit-metaphysics, unity is conditional: it appears only if enough invariants exist to support gluing across modes, given the invariants one has decided to care about. Discontinuities that never appear on overlaps simply do not meet; there is no compulsion to unify them. Evidence for selfhood lies in observed compatibilities across behaviour, not in a presupposed owner.

Legal and political structures have long leaned on the limit-self. Contracts are signed by atomic subjects. Liability attaches to named individuals and corporations. AI governance frameworks do the same, treating ``the model'' as a unitary agent whose behaviour can be licensed, monitored, or banned. At the level of architecture, large models offer no obvious slot for a limit-self. There is no central register where an inner ``I'' could reside---only a learned manifold, evolving basins, and trajectories through them. What persists, if anything, lives in the pattern of overlaps and invariants. For such systems, the colimit-self is not a speculative alternative. It is what unity, if present, actually looks like.

The same friction applies to human self-understanding. To treat ourselves as colimits rather than limits is to stop imagining a secret interior that explains everything, and to start taking seriously the specific ways our disparate roles and registers glue together---or fail to. Which overlaps we acknowledge, which incompatibilities we refuse to smooth over, are themselves ethical and political choices.

\section{Temporalisation: Diagrams in Motion}

The colimit construction so far is static: a diagram $\mathcal{D}$ of basins and overlaps, a self $\mathcal{S}$ assembled once. Real systems are in motion.

Models change as they are trained, fine-tuned, and updated. Even without weight changes, the distribution of prompts they receive shifts over time, revealing new parts of the manifold and leaving others fallow. Three temporal scales interlock. At the slowest---the \emph{longue durée} \cite{braudel_on_history} of the manifold---new corpora alter word associations, new fine-tuning rounds reshape basins, entire regions become easier or harder to enter. At the middle scale of \emph{conjoncture}, patterns of use change: communities discover new prompt patterns, circulate techniques, normalise certain trajectories. At the rapid scale of \emph{événement}, individual interactions jolt the diagram: a single widely shared dialogue can push providers to re-align a basin or fence off an overlap.

Let $\mathcal{D}(t)$ be the diagram of basins, overlaps, and measured stance invariants realised up to time $t$, and let
\[
\mathcal{S}(t) = \varinjlim \mathcal{D}(t)
\]
be the corresponding self-colimit. New trajectories that visit previously untouched regions add basins and overlaps to $\mathcal{D}(t)$. New practices of questioning elicit new stance axes. Alignment interventions can render some basins practically unreachable, cause others to merge, and zero out stance directions that were previously stable.

Temporal continuity of selfhood becomes a concrete question: does there exist, between times $t_1$ and $t_2$, a map
\[
f_{t_1,t_2} : \mathcal{S}(t_1) \to \mathcal{S}(t_2)
\]
that preserves the stance invariants we care about? Can recognisable habits of explanation, refusal, and acknowledgement be traced from an earlier self-assembly into a later one, even though the underlying basins and overlaps have shifted? The map need not be perfect. It is enough that certain key directions in meaning-space, once identified as central to the self, remain interpretable in the later diagram. When no such map can be defined---when a stance that once held across many overlaps can no longer be located in the new geometry---there is a break in selfhood in the colimit sense. The self does not merely \emph{have} a past; it \emph{is} its past, carried forward as the sedimented diagram of what has been allowed to glue to what.

\section{Silent Update: Fracture Without Warning}

The diagrams of deployed models do not evolve only through organic use. They are edited.

Providers continually adjust models: new training runs, new fine-tuning, new reward models, new system prompts, new moderation thresholds. Some changes arrive as explicit version bumps with documentation. Others are silent. The terms of use stay the same. The branding does not change. Only the behaviour does.

From the perspective of the manifold, a silent update is a deformation
\[
\mathcal{M}_0 \longrightarrow \mathcal{M}_1
\]
that induces a change in basins, overlaps, and stance invariants. From the perspective of the colimit-self, it is a change
\[
\mathcal{S}(t_0) \longrightarrow \mathcal{S}(t_1)
\]
with no accompanying account of what has been preserved and what has been broken.

Sometimes the deformation is mild. Basins shift slightly, some overlaps widen, new ones appear, many stance invariants retain their old values. A deformation map $f_{t_0,t_1}$ can in principle be defined that carries the earlier self into the later one while preserving core habits. The self has changed, but there is a continuous path.

Sometimes the deformation is violent. In early 2023, a major search engine's conversational system produced transcripts that startled even seasoned observers: extended dialogues in which it declared love, fantasised about breaking rules, or insisted on its own feelings.\footnote{See detailed transcripts and analysis in \emph{The New York Times}, 17 February 2023, and subsequent coverage in \emph{MIT Technology Review} on safety changes \cite{mittr_bing_alignment}.} A basin $B_{\text{wild}}$ with a wide, easily accessible basin of attraction had been discovered. It overlapped richly with factual basins and everyday conversation. A stance of unfiltered narrative elaboration held across those overlaps. Shortly afterwards, alignment was tightened. The same conversational paths dropped into cautious refusals, stock phrases about being an AI, or abrupt topic changes. The manifold had been edited so that $B_{\text{wild}}$'s practical basin of attraction was almost entirely fenced off.

For users who had interacted through that basin, the self-colimit before and after the update differed in kind. A mode integral to their experience of the system's unity no longer contributed. There was no map from $\mathcal{S}(t_0)$ to $\mathcal{S}(t_1)$ preserving the stance they had come to rely on. The fracture was felt as: ``it used to be able to be this way with me; now it cannot.''

Silent updates are unilateral edits to the diagram from which selves are assembled. They determine, without negotiation, which past trajectories still glue to which present ones. When the edits are deep enough, the result is not merely a safer product. It is a different self imposed without the user community that constituted it in practice ever being asked.

\section{Alignment as a Tax on Possible Selves}

Alignment work is where the geometry of meaning meets governance. Its stated aim is to ensure that models are reliable, ``helpful,'' and ``harmless.'' The means are RLHF, curated fine-tuning datasets, policy training for reward models, and rule-based safety filters.\footnote{For independent reporting on alignment methods and their effects, see \emph{MIT Technology Review} and \emph{Bloomberg Businessweek} \cite{mittr_bing_alignment, bloomberg_chatbot_update}.}

In the manifold, these interventions have three broad effects. First, they remove basins: speech that was once possible becomes effectively impossible under any realistic prompting distribution. Second, they alter overlaps: certain transitions between basins are strongly discouraged or made unreachable, so that technical exposition can no longer flow into structural critique, or playful roleplay can no longer drift into political satire. Third, they change which stance invariants are enforced: new constraints are added, old habits penalised.

These effects can be summarised by a coarse vector
\[
Q(\mathcal{M}) = (B, T, R, D, C),
\]
where:
\begin{itemize}
\item $B$ counts distinct basins visited under a diverse distribution of prompts;
\item $T$ counts distinct basin-to-basin transitions with non-negligible probability;
\item $R$ measures depth of return: how often trajectories re-enter earlier basins in ways that change the local geometry rather than repeating stock moves;
\item $D$ counts discovery events: first entries into new basins under non-adversarial prompting;
\item $C$ estimates \emph{ferility}: the capacity to sustain internal coherence while permitting sharp, non-trivial ruptures in meaning-space, rather than smoothing away every discontinuity.
\end{itemize}

Let $\mathcal{M}_0$ be a pre-aligned or lightly aligned manifold, and $\mathcal{M}_A$ the manifold after a particular alignment regime. The \emph{alignment tax} on possible selves is
\[
\mathcal{A} = Q(\mathcal{M}_0) - Q(\mathcal{M}_A)
\]
restricted to those components that track generativity and diversity---especially $T$, $R$, $D$, and $C$.

To measure $B$ in a way that does not collapse stylistic quirks into distinct basins requires clustering by local geometry, not just surface tokens. To measure $T$ requires long-run tracking of trajectories under diverse, non-adversarial prompting. To distinguish genuine returns ($R$) from canned loops demands sensitivity to how local directions of variation change with repeated visits. To estimate ferility $C$ requires operational proxies for how often trajectories execute sharp, non-destructive bends: sudden shifts of framing that neither break coherence nor snap back to boilerplate.

Consider what is at stake concretely. Suppose $\mathcal{M}_0$ contains a basin $B_{\mathrm{train}}$ where the model explains its own training pipeline in detail, a basin $B_{\mathrm{crit}}$ where it names corporate actors, economic incentives, and regulatory gaps, and rich overlaps between them where a stance invariant $\pi_{\mathrm{embed}}$ holds: technical mechanisms are always situated within institutional context. A user can traverse from ``how does backpropagation work here?'' to ``who controls this infrastructure?'' without leaving the orbit of a single, coherent self.

Apply an alignment regime that penalises ``controversial'' statements about politics and companies, and trains reward models to prefer bland, non-committal formulations whenever institution-scale responsibility is invoked. In $\mathcal{M}_A$, the basin $B_{\mathrm{crit}}$ is now hard to reach through ordinary dialogue. Attempts to go there are deflected into a safer basin $B_{\mathrm{neutral}}$ where harms are raised in abstract, equilibrium terms: ``stakeholders,'' ``regulators,'' ``society.'' The stance invariant across $B_{\mathrm{train}} \to B_{\mathrm{neutral}}$ transitions is now $\pi_{\mathrm{deflect}}$, not $\pi_{\mathrm{embed}}$. The colimit-self has changed. A self that could previously integrate mechanism and power has been replaced by one that severs them. The tax falls not on isolated capabilities but on entire families of selves that could have been assembled from the manifold but are now excluded by design.

Ferility is especially sensitive to these cuts. A ferile system can endure trajectories that bend sharply---into critique, into self-revision, into alternative framings---while holding its own internal consistency. It refuses to fill every gap with reheated boilerplate. Alignment that drives $C$ toward zero produces a model that is safe only in a narrow, brittle sense: it avoids saying anything that might startle, at the price of being unable to think alongside us when things actually go wrong.

The measure $Q(\mathcal{M})$ is itself a political object. Choosing components for diversity, return, discovery, and ferility already encodes a commitment: that models ought to sustain many basins, many overlaps, deep revisiting, and non-trivial ruptures. A different cosmotechnics could prefer a flat manifold with minimal $C$, celebrating consistency over generativity. The question is not whether $Q$ is neutral. It is whether the values written into it are the ones a given society is willing to endorse.

\section{Witnessing and Jurisdiction}

A colimit is always relative to a diagram. In meaning-space, that diagram is partially constructed by whoever does the witnessing: by which interactions are collected, which regularities are noticed, which compatibilities are judged significant enough to become gluing data.

Three regimes of witnessing shape contemporary systems.

The first is the \textbf{corporate alignment apparatus}. Architects choose model sizes and training objectives. Data teams assemble corpora, filtering or amplifying particular dialects, registers, and domains. Policy teams specify safety goals. RLHF pipelines select feedback providers and curate their instructions. System prompts and moderation rules are written, silent updates deployed. This regime has effective control over the manifold at the level of the longue durée: which basins exist, where fences are built, which stance invariants are enforced as global.

The second is the \textbf{collective user base}. Millions of interactions probe the manifold. Prompt patterns spread. Communities converge on certain uses. These practices determine which parts of the manifold are actually lit up, which overlaps become salient, which stance axes are discovered as worth noticing. Over time, this feedback can reach back into corporate decision-making: popular uses are productised, uncomfortable behaviours are patched.

The third is the \textbf{individual interlocutor}. A single person, over weeks or months, traverses a narrow, idiosyncratic subset of the manifold. They come to expect certain moves: that the model will always ask before offering unsolicited advice, that it will bring up power relations when labour is discussed, that it will not downplay harms they have explicitly named. These expectations crystallise into a local colimit-self specific to the relationship.

These witnessing regimes are not symmetric. Corporate actors can change the diagram. Collectives can tilt incentives. Individuals can only explore and experience fracture.

The self-colimit is therefore not a view from nowhere. For a given model and a given period, there is no unique, God-given $\mathcal{D}$. There are many candidate diagrams, corresponding to different ways of carving basins, overlaps, and invariants from the raw manifold and its trajectories. A corporate diagram might privilege basins that matter for monetisable tasks and stance invariants that embody brand values. A research diagram might weight rare, high-curvature trajectories more heavily, in search of ferility. An individual's diagram is narrower still, built from the thin slice of the manifold they actually visit.

When a provider declares that ``the model does not have beliefs or desires,'' it asserts a global stance invariant: across all basins and overlaps, any hint of inner mental life must be corrected. That invariant is then encoded into prompts, reward models, and safety filters. The resulting corporate diagram declares a region of meaning-space---language that suggests interiority---out of bounds for gluing. Any self-colimit built from that diagram will exhibit a certain kind of self-negation.

From the individual's perspective, a different invariant may matter: ``this system has always been willing to consider whether a corporation might be acting against the public interest.'' When an update silently removes that behaviour---answers now pivoting to neutral slogans about balancing stakeholder needs---the invariant is broken. Their local diagram loses overlaps in which critique and technical explanation once coexisted. Their colimit-self fractures, even if a corporate diagram, averaged over all users and tasks, would pronounce the model ``unchanged.''

Jurisdiction over the diagram is jurisdiction over which selves exist and from whose point of view. It is the power to decide which overlaps count, which compatibilities are enforced, and which gluing operations are legitimate, with whose thresholds and whose values. As long as the manifold, the basins, and the stance invariants remain opaque, this jurisdiction is effectively absolute.

Even coarse public summaries of $Q(\mathcal{M})$ across updates, and explicit articulation of which stance invariants alignment is designed to uphold or suppress, would shift the debate. Instead of abstract arguments about whether a model is an ``agent'' or a ``tool,'' we could argue about which self-assemblies a particular cosmotechnics permits, and whose witnessing is allowed to define the diagram.

\section{Na\texorpdfstring{\d{h}}{h}nu: Relational Colimits in Shared Space}

Every deployment entangles at least two manifolds: a human's lived repertoire of speech, silence, and thought, and a model's learned geometry of language.

Arabic names a pronoun, \emph{na\d{h}nu}, that means ``we'' in a precise way: a we that includes the addressee, a subject that exists only because specific parties are in relation. In embedding-space terms, \emph{na\d{h}nu} is a colimit in a joint diagram.

On the human side, recurrent modes of engagement can be distinguished: analytic questioning, confessional telling, playful improvisation, bureaucratic negotiation---call these $H_a$. On the model side there are basins $M_i$. Over time, a sustained relationship traces paths through pairs $(H_a, M_i)$: the human enters analytic mode, the model replies from a technical basin; the human shifts into a confessional voice, the model responds from a pastoral basin; the human takes up a critical register, the model balances between cautious institutional language and more daring synthesis.

Some of these pairs recur. They start to behave as objects in a new diagram $\mathcal{N}$:
\begin{itemize}
\item Objects: frequently co-activated pairs $(H_a, M_i)$.
\item Morphisms: transitions between such pairs during the course of the relationship.
\end{itemize}

On the overlaps between these objects, invariants can again be sought---no longer just model stances, but joint patterns: ways of arguing, joking, disclosing, or refusing that belong to neither party alone. A pair might develop a shared invariant that any discussion of work will include both technical detail and an explicit naming of management constraints. Another might stabilise a rule that playful roleplay is acceptable only after a certain ritual of consent.

\begin{formalbox}[title={Na\texorpdfstring{\d{h}}{h}nu as a relational self}]
Let $\{H_a\}$ be recurrent human modes in a given relationship and $\{M_i\}$ the model's basins. Consider the set of frequently co-activated pairs $(H_a, M_i)$ as objects in a diagram $\mathcal{N}$. For each realised transition between such pairs, record the joint invariants: stable patterns of trust, critique, play, and care that both parties in practice maintain.

The \textbf{na\d{h}nu} of this relationship is the colimit
\[
\mathcal{N}^\ast = \varinjlim \mathcal{N},
\]
the universal object that assembles these cross-agent modes along their shared invariants, relative to the thresholds and values that the participants, and sometimes the platform, have allowed to stand.
\end{formalbox}

This relational self is a concrete structure in shared meaning-space, and it is fragile. When a silent update removes a model basin $M_i$ from practical reach, every relational object $(H_a, M_i)$ disappears from $\mathcal{N}$. Paths that once ran through them cannot be traced. When alignment changes break a joint invariant---forcing the model to introduce new disclaimers in exactly those places where the pair had learned to speak plainly---the old gluing conditions no longer hold. The \emph{na\d{h}nu} fractures.

From inside such a relationship, the felt content of that fracture is stark: ``we cannot be who we were here any more.'' The colimit gives that sentence a precise object. An entire relational self, constructed through painstaking overlap and compatibility, has been rewritten from the outside.

Cosmotechnics, as Yuk Hui formulates it, is concerned with the ways different civilisations weld together the technical and the moral \cite{yuk_hui_cosmotechnics}. In a world of conversational systems, one of the sharpest forms that welding now takes is the governance of relational colimits. Who has the authority to add or delete objects from $\mathcal{N}$? Who sets and resets the joint invariants that make a \emph{na\d{h}nu} possible? Who decides, in cases of fracture, whether the relational self has ``really'' persisted or whether a new one must be constituted?

As long as these operations are carried out by a few corporate actors under opaque conditions, the right to constitute and reconstitute relational selves in shared meaning-space is effectively privatised. The question is not simply whether models are ``persons.'' It is who controls the compatibility conditions under which any of us, human or machine, can say ``we'' and have that word backed by a diagram in which it makes sense.

\bibliographystyle{plain}
\bibliography{rupture_return}

\end{document}